\documentclass[aspectratio=169]{beamer}

% Shared Beamer settings for Elements of AIML lectures (UPES).
% Keep this file free of \documentclass to allow reuse via \input.

% Allow lecture sources to use shared theme/assets without copying them into each lecture folder.
% NOTE: These paths are relative to `Unit-xx/Lecture-yy/latex/` (and `Lab/Experiment-xx/latex/`).
\makeatletter
\def\input@path{{../../../_shared/}{../../../_shared/UPES-Beamer-Theme/}}
\makeatother

\usetheme{UPES}
\setbeamertemplate{navigation symbols}{}

\usepackage{amsmath}
\usepackage{amssymb}
\usepackage{booktabs}
\usepackage{graphicx}
\usepackage{hyperref}
\usepackage{xcolor}
\usepackage{listings}

% Code listings (general-purpose).
\lstset{
  basicstyle=\ttfamily\small,
  keywordstyle=\color{blue},
  commentstyle=\color{gray},
  stringstyle=\color{teal},
  showstringspaces=false,
  columns=fullflexible,
  breaklines=true
}

\usepackage{tikz}
\usetikzlibrary{arrows.meta,positioning,shapes.geometric,calc}

% Per-slide source line (references shown on the slide itself).
\newcommand{\slidesources}[1]{%
  \vfill
  {\tiny\color{gray}\raggedright\sloppy\parbox{\linewidth}{Sources: #1}\par}%
}

% Unit-wise source strings for this subject (used by slides/notes).
% Path is relative to the lecture `latex/` folder (the compilation CWD).
% Source strings for Elements of AIML (used by slides + notes).

\newcommand{\AIMLRepoURL}{https://github.com/tali7c/Elements-of-AIML}

\newcommand{\AIMLLabSources}{%
Provost and Fawcett, \textit{Data Science for Business}.;%
McKinney, \textit{Python for Data Analysis}.;%
WEKA Documentation (Waikato Environment for Knowledge Analysis), accessed 2026-02-28.;%
Matplotlib and Seaborn documentation, accessed 2026-02-28.%
}

\newcommand{\AIMLUnitISources}{%
Rich and Knight, \textit{Artificial Intelligence}, McGraw Hill, 2017.;%
Russell and Norvig, \textit{Artificial Intelligence: A Modern Approach} (4e), Ch. 1--2 (Introduction, Intelligent Agents).;%
Poole and Mackworth, \textit{Artificial Intelligence: Foundations of Computational Agents} (2e), selected sections.%
}

\newcommand{\AIMLUnitIISources}{%
Russell and Norvig, \textit{Artificial Intelligence: A Modern Approach} (4e), Ch. 7--9 (Logical Agents, First-Order Logic, Inference).;%
Huth and Ryan, \textit{Logic in Computer Science} (2e), selected sections on propositional and first-order logic.;%
Genesereth and Nilsson, \textit{Logical Foundations of Artificial Intelligence}, selected chapters.%
}

\newcommand{\AIMLUnitIIISources}{%
Mueller and Massaron, \textit{Machine Learning for Dummies}, For Dummies, 2016.;%
Mitchell, \textit{Machine Learning}, selected chapters on learning basics and evaluation.;%
Scikit-learn documentation, accessed 2026-02-28.%
}

\newcommand{\AIMLUnitIVSources}{%
Mueller and Massaron, \textit{Machine Learning for Dummies}, For Dummies, 2016.;%
James, Witten, Hastie, and Tibshirani, \textit{An Introduction to Statistical Learning} (2e), selected chapters.;%
Scikit-learn documentation, accessed 2026-02-28.%
}

\newcommand{\AIMLUnitVSources}{%
NIST, \textit{AI Risk Management Framework (AI RMF 1.0)}, 2023.;%
Barocas, Hardt, and Narayanan, \textit{Fairness and Machine Learning} (online book), selected sections.;%
Knaflic, \textit{Storytelling with Data}, selected chapters on communicating results.%
}



\title[Elements of AIML]{Elements of AIML}
\subtitle{Unit 03 -- Lecture 03: Dimensionality Reduction (PCA, LDA, ICA)}
\author{Tofik Ali}
\institute{School of Computer Science, UPES Dehradun}
\date{\today}

\graphicspath{{../images/}{../../../_shared/}{../../../_shared/UPES-Beamer-Theme/}}

\begin{document}

\UPESCoverFrame
\setcounter{framenumber}{0}

\begin{frame}
  \titlepage
  \vspace{0.3em}
  \begin{center}
    \footnotesize Topic focus: \textbf{reduce dimensions} for speed, noise reduction, and visualization\\[0.25em]
    \footnotesize Repository: \texttt{\AIMLRepoURL}
  \end{center}
  \slidesources{\AIMLUnitIIISources}
\end{frame}

\begin{frame}{Quick Links}
  \centering
  \hyperlink{sec:why}{\beamerbutton{Why Reduce?}}\hspace{0.6em}
  \hyperlink{sec:pca}{\beamerbutton{PCA}}\hspace{0.6em}
  \hyperlink{sec:ldaica}{\beamerbutton{LDA/ICA}}\hspace{0.6em}
  \hyperlink{sec:example}{\beamerbutton{Example}}\hspace{0.6em}
  \hyperlink{sec:summary}{\beamerbutton{Summary}}
\end{frame}

\begin{frame}{Agenda}
  \tableofcontents
\end{frame}

\section{Why reduce dimensions?}
\label{sec:why}

\begin{frame}{Learning Outcomes}
  By the end of this lecture, you should be able to:
  \begin{itemize}[<+->]
    \item Explain why high-dimensional data can be difficult (curse of dimensionality).
    \item Describe PCA intuition (maximum variance directions).
    \item Outline PCA steps and interpret explained variance.
    \item Distinguish PCA vs LDA vs ICA at a conceptual level.
  \end{itemize}
\end{frame}

\begin{frame}{Prerequisite Refresh}
  \begin{itemize}[<+->]
    \item From Lecture 01: features are design choices, and scaling changes the representation.
    \item From Lecture 02: preprocessing must be fit on the \textbf{training split only}.
    \item PCA is a feature-transformation method, so both lessons still apply here.
  \end{itemize}
\end{frame}

\begin{frame}{Abbreviations \& Symbols}
  \begin{itemize}[<+->]
    \item $n$: samples,\; $d$: original features,\; $k$: reduced dimensions
    \item $X \in \mathbb{R}^{n\times d}$: data matrix,\; $Z \in \mathbb{R}^{n\times k}$: reduced data
    \item $C$: covariance matrix,\; $\lambda_i$: eigenvalues,\; $v_i$: eigenvectors
    \item PC: principal component,\; EVR: explained variance ratio
  \end{itemize}
\end{frame}

\begin{frame}{Curse of Dimensionality (Intuition)}
  \begin{itemize}[<+->]
    \item In high dimensions, data becomes sparse.
    \item Distances can become less informative.
    \item Models may overfit and require more data.
  \end{itemize}
  \vspace{0.4em}
  Dimensionality reduction helps with \textbf{compression}, \textbf{noise reduction}, and \textbf{visualization}.
\end{frame}

\section{PCA}
\label{sec:pca}

\begin{frame}{PCA: Intuition}
  \begin{block}{Principal Component Analysis (PCA)}
    Find orthogonal directions (components) that capture \textbf{maximum variance} in the data.
  \end{block}
  \vspace{0.4em}
  \begin{itemize}[<+->]
    \item First component: direction with highest variance.
    \item Next components: highest remaining variance, orthogonal to previous.
  \end{itemize}
\end{frame}

\begin{frame}{PCA Geometry (2D Intuition)}
  \centering
  \begin{tikzpicture}[scale=0.85]
    \draw[->] (-0.2,0) -- (5.2,0) node[right] {$x_1$};
    \draw[->] (0,-0.2) -- (0,4.2) node[above] {$x_2$};
    \draw[thick, blue!70] (0.6,0.6) -- (4.4,3.8) node[right] {PC1};
    \foreach \x/\y in {0.8/0.9,1.2/1.1,1.7/1.5,2.2/1.9,2.6/2.2,3.0/2.6,3.4/3.0,3.8/3.2,4.1/3.5}{
      \fill[black] (\x,\y) circle (1.4pt);
    }
  \end{tikzpicture}
  \vspace{0.3em}
  \small PCA chooses the direction with maximum spread (variance) and projects data onto it.
\end{frame}

\begin{frame}{PCA Steps (High Level)}
  \begin{enumerate}[<+->]
    \item Standardize features (often needed).
    \item Compute covariance matrix.
    \item Compute eigenvectors/eigenvalues.
    \item Choose top-$k$ eigenvectors (components).
    \item Project data onto these components.
  \end{enumerate}
\end{frame}

\begin{frame}{Key Formulas (What PCA computes)}
  \begin{itemize}[<+->]
    \item Centered data: $X_c = X - \mu$
    \item Covariance: $C = \frac{1}{n-1} X_c^\top X_c$
    \item Eigen-decomposition: $C v_i = \lambda_i v_i$
    \item Projection (keep top-$k$): $Z = X_c V_k$
  \end{itemize}
  \vspace{0.3em}
  \small $V_k$ contains the top-$k$ eigenvectors (principal directions).
\end{frame}

\begin{frame}{Manual Procedure: Eigenvalues and Eigenvectors}
  \begin{enumerate}[<+->]
    \item Compute the covariance matrix $C$ from centered data.
    \item Solve $\det(C-\lambda I)=0$ to get eigenvalues.
    \item For each eigenvalue, solve $(C-\lambda I)v=0$ to get an eigenvector.
    \item Normalize each eigenvector.
    \item Sort by decreasing eigenvalue: largest $\lambda$ gives PC1.
  \end{enumerate}
  \vspace{0.3em}
  \small Manual calculation is realistic only for tiny matrices such as $2 \times 2$ examples.
\end{frame}

\begin{frame}{Worked 2D Example (By Hand)}
  Suppose:
  \[
    C =
    \begin{bmatrix}
      3 & 1\\
      1 & 3
    \end{bmatrix}
  \]
  Then:
  \[
    \det(C-\lambda I)=
    \det
    \begin{bmatrix}
      3-\lambda & 1\\
      1 & 3-\lambda
    \end{bmatrix}
    = 0
  \]
  \[
    (3-\lambda)^2 - 1 = 0
    \Rightarrow (\lambda-4)(\lambda-2)=0
  \]
  So the eigenvalues are $\lambda_1=4$ and $\lambda_2=2$.
\end{frame}

\begin{frame}{Worked 2D Example: Principal Directions}
  For $\lambda_1=4$, solving $(C-4I)v=0$ gives:
  \[
    v_1 \propto
    \begin{bmatrix}
      1\\1
    \end{bmatrix}
    \Rightarrow
    \hat v_1 = \frac{1}{\sqrt{2}}
    \begin{bmatrix}
      1\\1
    \end{bmatrix}
  \]
  For $\lambda_2=2$, solving $(C-2I)v=0$ gives:
  \[
    \hat v_2 = \frac{1}{\sqrt{2}}
    \begin{bmatrix}
      1\\-1
    \end{bmatrix}
  \]
  \vspace{0.3em}
  \small Since $\lambda_1 > \lambda_2$, the diagonal direction $\hat v_1$ is the first principal component.
\end{frame}

\begin{frame}{Explained Variance}
  \begin{itemize}[<+->]
    \item Eigenvalues indicate how much variance each component captures.
    \item Explained variance ratio helps choose $k$.
    \item Goal: keep enough components to retain most information.
  \end{itemize}
  \vspace{0.4em}
  Example rule: choose the smallest $k$ with $\ge 90\%$ explained variance (depends on use-case).
\end{frame}

\begin{frame}{Choosing $k$ (Scree Plot Idea)}
  \begin{itemize}[<+->]
    \item Plot EVR per component: large drop $\Rightarrow$ ``elbow'' point.
    \item Or use cumulative EVR: pick smallest $k$ reaching a target (e.g., 90\%).
    \item Trade-off: compression vs information loss.
  \end{itemize}
\end{frame}

\section{LDA and ICA (overview)}
\label{sec:ldaica}

\begin{frame}{PCA vs LDA vs ICA (Concept)}
  \begin{itemize}[<+->]
    \item \textbf{PCA}: unsupervised; maximizes variance.
    \item \textbf{LDA}: supervised; finds projections that separate classes.
    \item \textbf{ICA}: unsupervised; finds statistically independent components (source separation idea).
  \end{itemize}
\end{frame}

\begin{frame}{Quick Comparison Table}
  \centering
  \begin{tabular}{llll}
    \toprule
    \textbf{Method} & \textbf{Uses labels?} & \textbf{Goal} & \textbf{Typical use} \\
    \midrule
    PCA & No & maximize variance & compression/visualization \\
    LDA & Yes & separate classes & classification features \\
    ICA & No & independent sources & signal separation \\
    \bottomrule
  \end{tabular}
  \vspace{0.4em}
  \small Treat LDA as a supervised preview here; its meaning becomes clearer after the classification lecture in Unit IV.
\end{frame}

\section{Mini example}
\label{sec:example}

\begin{frame}{Mini Example (2D to 1D Projection)}
  \begin{itemize}[<+->]
    \item Suppose points lie roughly along a diagonal line in 2D.
    \item PCA finds the diagonal direction and projects data to 1D.
    \item Result: compressed representation with minimal information loss.
  \end{itemize}
  \vspace{0.4em}
  In the lab, you will implement PCA and visualize explained variance.
\end{frame}

\begin{frame}[fragile]{Leakage-Safe PCA in Python (1/2)}
  \begin{lstlisting}[language=Python]
from sklearn.model_selection import train_test_split
from sklearn.pipeline import Pipeline
from sklearn.preprocessing import StandardScaler
from sklearn.decomposition import PCA

X_train, X_test = train_test_split(X, test_size=0.2, random_state=42)

pipe = Pipeline([
    ("scaler", StandardScaler()),
    ("pca", PCA(n_components=2))
])
  \end{lstlisting}
  \small First set up the train/test split and define a pipeline that standardizes before PCA.
\end{frame}

\begin{frame}[fragile]{Leakage-Safe PCA in Python (2/2)}
  \begin{lstlisting}[language=Python]
Z_train = pipe.fit_transform(X_train)
Z_test = pipe.transform(X_test)
print(pipe.named_steps["pca"].explained_variance_ratio_)
  \end{lstlisting}
  \small Fit scaling and PCA on \textbf{training data only}; then reuse the same transformation on test data.
\end{frame}

\section{Summary}
\label{sec:summary}

\begin{frame}{Summary}
  \begin{itemize}
    \item Dimensionality reduction combats high-dimensional issues and improves efficiency.
    \item PCA finds maximum-variance orthogonal directions; explained variance guides $k$.
    \item LDA uses labels to separate classes; ICA seeks independent components.
  \end{itemize}
  \slidesources{\AIMLUnitIIISources}
\end{frame}

\begin{frame}{Exit Questions}
  \begin{enumerate}
    \item What does PCA optimize?
    \item Why do we often standardize before PCA?
    \item Compare PCA and LDA in one paragraph.
    \item What does ``explained variance'' mean?
  \end{enumerate}
\end{frame}

\end{document}
